Function calling enables Large Language Models (LLMs) to interact with external tools and APIs by translating natural-language requests into structured function invocations. Given a user request and a set of available functions, an LLM must decide whether to call a function, select the appropriate one, and generate arguments that reflect the user's intent. This capability is central to LLM-based agentic systems that retrieve information, execute actions, and make decisions in external environments. However, syntactic validity does not guarantee semantic correctness: a call can satisfy a schema while selecting the wrong function or filling arguments with incorrect values.

Recent work improves function-calling reliability through schema validation, tool-use verification, and iterative refinement. Schema-based methods enforce structural constraints, while self-refinement uses model-generated critiques to identify and correct errors. Complementary verification-oriented methods explicitly assess function selection and argument values, and feedback-driven agent frameworks revise subsequent actions using signals such as tool-use error detection or execution observations.

Building on these lines of work, recent advances in decision modeling offer a promising approach to providing targeted judgments over candidate outputs. Unlike conventional LLM-based critiques, decision models can produce structured binary judgments to assess whether specific evaluation criteria are satisfied. Recent approaches further explore confidence-aware verification, where lightweight decision models evaluate candidate outputs and selectively defer uncertain cases to stronger judges, improving evaluation efficiency. 

\textcolor{red}{However, turning such judgments into reliable improvements remains challenging. \textbf{First}, structural validation alone cannot determine whether a function call aligns with the user's request, particularly when the correct action is to call no function. \textbf{Second}, function selection and argument generation are distinct error sources: a model may choose the right function but provide wrong arguments, or generate plausible arguments for the wrong function, which calls for fine-grained verification. \textbf{Third}, although iterative refinement can correct mistakes, using only the latest feedback can discard useful signals from earlier failed attempts. These challenges motivate leveraging decision-model verification during refinement while retaining verification history across attempts.
}

To this end, we propose \textbf{Decision-Model-Guided Function Calling (DMFC)}, an inference-time framework that integrates decision-model-based verification with cumulative feedback-driven refinement. \textcolor{red}{DMFC evaluates each generated call with two complementary binary judgments---function selection correctness and argument correctness---in addition to structural schema validation. When a candidate fails verification, the rejected call and its judgments are stored as cumulative feedback and injected into the LLM context to guide subsequent revisions.} A call is accepted only when all applicable verification criteria are satisfied, enabling iterative correction without executing external functions.

\textcolor{red}{We evaluate DMFC using diverse decision-model verifier on the Berkeley Function Calling Leaderboard (BFCL) and xxx.} DMFC improves accuracy from 64.1\% to 74.8\% with GPT-5-nano and from 81.1\% to 82.1\% with GPT-5.4-mini. Retaining verification history further outperforms latest-only feedback by 1.1 percentage points, with the largest gains from preventing unnecessary function calls at the cost of additional inference compute.

Our main contributions are threefold:
\textcolor{red}{
\begin{itemize}
    \item We introduce a decision-model-guided verification framework that explicitly evaluates function selection and argument correctness before execution.
    \item We propose a cumulative feedback-driven refinement mechanism that leverages historical verification judgments to improve subsequent function calls.
    \item We conduct comprehensive experiments and ablation studies to evaluate the effectiveness of the proposed framework in terms of function-calling accuracy, correction behavior, and computational efficiency.
\end{itemize}}